\documentclass[11pt]{article}
\usepackage[margin=1in]{geometry}
\usepackage{amsmath,amssymb,amsthm,mathtools}
\usepackage[hidelinks]{hyperref}

\newtheorem{theorem}{Theorem}
\newtheorem{proposition}[theorem]{Proposition}
\newtheorem{lemma}[theorem]{Lemma}
\newtheorem{corollary}[theorem]{Corollary}
\theoremstyle{remark}
\newtheorem{remark}[theorem]{Remark}

\newcommand{\R}{\mathbb R}
\newcommand{\one}{\mathbf 1}
\newcommand{\tr}{\operatorname{tr}}
\newcommand{\supp}{\operatorname{supp}}
\newcommand{\ip}[2]{\langle #1,#2\rangle}

\title{Equality in the Positive Square-Energy Bound}
\author{}
\date{25 July 2026}

\begin{document}
\maketitle

\begin{abstract}
Let $s^+(G)$ be the sum of the squares of the positive adjacency
eigenvalues of a graph $G$.  Liu, Tang, and Zhang recently proved the
Elphick--Farber--Goldberg--Wocjan lower bound $s^+(G)\geq |V(G)|-1$ for
connected graphs.  We determine its equality case: a connected graph $G$ of
order $n$ satisfies $s^+(G)=n-1$ if and only if $G$ is a tree.  We first
trace equality through the doubly nonnegative inequality of Liu--Tang--Zhang.
Their folding and cut-vertex induction, supplemented by their averaged
vertex-deletion estimate, forces every block of an equality graph to be
complete.  We then prove that the positive spectral part of a tree is
strictly positive between vertices having a common neighbor.  A spanning-tree
variational argument excludes every non-tree block graph.
\end{abstract}

\section{Introduction}

All graphs are finite, simple, and undirected.  Let $A=A(G)$ have eigenvalues
$\lambda_1,\ldots,\lambda_n$, and define
\[
 s^+(G)=\sum_{\lambda_i>0}\lambda_i^2,
 \qquad
 s^-(G)=\sum_{\lambda_i<0}\lambda_i^2.
\]
Writing the positive and negative spectral parts of $A$ as
\[
 A=P-B,\qquad P=A_+\succeq0,\quad B=A_-\succeq0,\quad PB=0,
\]
we have $s^+(G)=\tr P^2$, $s^-(G)=\tr B^2$, and
\begin{equation}\label{eq:sum-energy}
 s^+(G)+s^-(G)=\tr A^2=2|E(G)|.
\end{equation}

Elphick, Farber, Goldberg, and Wocjan conjectured that
$\min\{s^+(G),s^-(G)\}\geq n-1$ for every connected graph~\cite{EFGW}.
Liu, Tang, and Zhang recently proved this theorem by means of a sharp
inequality over the doubly nonnegative cone~\cite{LTZ}.  Akbari, Kumar,
Mohar, Pragada, and Zhang asked for the equality cases and, in particular,
conjectured that equality for $s^+$ characterizes trees
\cite[Conjecture~9.2(i)]{AKMPZ}.  We prove that conjecture.

\begin{theorem}\label{thm:main}
Let $G$ be a connected graph on $n$ vertices.  Then
\[
 s^+(G)=n-1
 \quad\Longleftrightarrow\quad
 G\text{ is a tree}.
\]
\end{theorem}

The equality trace in Section~\ref{sec:equality} uses the ingredients of the
proof of the Liu--Tang--Zhang doubly nonnegative inequality, together with an
equality comparison not needed for their theorem.  We include the details,
both to fix the direction of the inequality and to make clear how the
complete-block conclusion follows.

\section{The doubly nonnegative inequality}\label{sec:dnn}

A real symmetric matrix $M$ is \emph{doubly nonnegative} if $M\succeq0$ and
$M$ is entrywise nonnegative.  For a connected graph $H$, put
\[
 q(H)=2|E(H)|-|V(H)|+1,
 \quad
 S(H,M)=\sum_{uv\in E(H)}\sqrt{M_{uv}},
 \quad
 T(M)=\one^TM\one.
\]
The orientation of the central inequality is important.

\begin{theorem}[Liu--Tang--Zhang \cite{LTZ}]\label{thm:ltz}
If $H$ is connected and $M$ is a doubly nonnegative matrix indexed by
$V(H)$, then
\begin{equation}\label{eq:ltz}
 4S(H,M)^2\leq q(H)T(M).
\end{equation}
\end{theorem}

We recall the parts of its proof used below.  First, every nonedge entry can
be folded onto the diagonal:
\begin{equation}\label{eq:fold}
 \mathcal F_H(M)=M+
 \sum_{\substack{\{u,v\}\notin E(H)\\u\ne v}}
 M_{uv}(e_u-e_v)(e_u-e_v)^T.
\end{equation}
The sum takes each unordered nonedge once.  The folded matrix is doubly
nonnegative, vanishes on nonedges, has the same edge entries as $M$, and has
the same value of $T$, since $(e_u-e_v)^T\one=0$.

Suppose next that an edge-supported pair $(H,M)$ is split at a cut vertex
into connected induced graphs $H_1,H_2$.  A Gram representation of $M$ and
orthogonal projection of the cut-vertex vector produce doubly nonnegative
matrices $M_1,M_2$ such that
\begin{equation}\label{eq:cut-add}
 q(H)=q(H_1)+q(H_2),\quad
 S(H,M)=S(H_1,M_1)+S(H_2,M_2),\quad
 T(M)=T(M_1)+T(M_2).
\end{equation}
The induction bounds for the two pairs are combined by scalar
Cauchy--Schwarz.

Finally, suppose $H$ has $h\geq3$ vertices and no cut vertex.  Then $H-v$ is
connected for every $v$.  For an edge-supported $M$, write
\[
 d_0=\tr M,
 \qquad w=\sum_{uv\in E(H)}M_{uv},
 \qquad T=d_0+2w,
 \qquad q=q(H).
\]
Averaging the induction hypothesis over the $h$ vertex deletions gives
\begin{equation}\label{eq:averaged}
 4S(H,M)^2\leq(q-1)T+\frac{q-1}{h-2}d_0.
\end{equation}
Independently, Cauchy--Schwarz on the $\ell=|E(H)|$ edges gives the flat
estimate
\begin{equation}\label{eq:flat}
 4S(H,M)^2\leq4\ell w.
\end{equation}
The proof of Theorem~\ref{thm:ltz} divides according as
\begin{equation}\label{eq:threshold}
 2(h-1)w\leq qd_0
\end{equation}
or the reverse strict inequality holds.  In the first case
\eqref{eq:flat} and $4\ell=2q+2(h-1)$ prove \eqref{eq:ltz}.  In the second
case, put $\beta=\ell-h+1$.  Then $q-1=(h-2)+2\beta$, and simplicity gives
\[
 q(h-2)-2\beta(h-1)=(h-1)(h-2)-2\beta\geq0.
\]
Multiplying the strict case hypothesis $qd_0<2(h-1)w$ by $h-2$ therefore
gives
\[
 2\beta(h-1)d_0\leq q(h-2)d_0
 <2(h-1)(h-2)w,
\]
and hence $\beta d_0<(h-2)w$, or equivalently
$(q-1)d_0/(h-2)<T$.  Thus \eqref{eq:averaged} gives the strict inequality
$4S^2<qT$.  These are precisely the folding, cut, averaged, and flat
ingredients from the proof in~\cite{LTZ}.

\section{Equality through folding and cuts}\label{sec:equality}

We now record the equality consequence needed later.  A \emph{block} is a
maximal subgraph with no cut vertex; thus a bridge is a block isomorphic to
$K_2$, while every block of order at least three is 2-connected.

\begin{proposition}[Equality trace]\label{prop:trace}
Let $H$ be a connected graph with at least one edge, let $M$ be doubly
nonnegative, and suppose
\begin{equation}\label{eq:tight}
 4S(H,M)^2=q(H)T(M)>0.
\end{equation}
Set $N=\mathcal F_H(M)$ and $\lambda=T(M)/q(H)$.  Then every block $C$ of
$H$ is either a bridge or a complete graph, and
\begin{equation}\label{eq:block-atoms}
 N=\lambda\sum_C\frac{q(C)}{|V(C)|^2}J_{V(C)},
\end{equation}
where each all-ones matrix is embedded on $V(C)$.  In particular, for every
edge $uv$ in a block $C$,
\begin{equation}\label{eq:edge-root}
 \sqrt{M_{uv}}=\sqrt{\lambda}\,
 \frac{q(C)}{2|E(C)|}.
\end{equation}
\end{proposition}

\begin{proof}
Folding preserves $S$ and $T$, so $(H,N)$ is again tight and is supported on
the diagonal and edges.  We trace the induction proving
Theorem~\ref{thm:ltz}.

At a cut vertex, use \eqref{eq:cut-add} and abbreviate
$q_i=q(H_i)$ and $T_i=T(N_i)$.  The induction and scalar
Cauchy--Schwarz give
\[
 2S(H,N)
 \leq \sqrt{q_1T_1}+\sqrt{q_2T_2}
 \leq \sqrt{(q_1+q_2)(T_1+T_2)}.
\]
Equality at the two ends forces equality in both descendant induction
bounds and equality in scalar Cauchy--Schwarz.  Hence
\begin{equation}\label{eq:ratio}
 \frac{T_1}{q_1}=\frac{T_2}{q_2}
 =\frac{T(N)}{q(H)}=\lambda.
\end{equation}
Repeating the split therefore leaves tight descendant pairs with the same
normalization $T_C=\lambda q(C)$.  The terminal descendants are bridges or
graphs of order at least three with no cut vertex.

We will also use the entrywise content of the cut construction.  Let $z$ be
the cut vertex, let $U$ and $W$ be the two sides away from $z$, and represent
$N$ by Gram vectors $g_x$.  Since $N$ is edge-supported, the spans
\[
 \mathcal U=\operatorname{span}\{g_u:u\in U\},\qquad
 \mathcal W=\operatorname{span}\{g_w:w\in W\}
\]
are orthogonal.  Decompose
\[
 g_z=p_U+p_W+p_0,
 \quad p_U\in\mathcal U,
 \quad p_W\in\mathcal W,
 \quad p_0\perp\mathcal U+\mathcal W.
\]
Use $p_U+p_0$ as the cut vector on the $U$-side and $p_W$ on the
$W$-side.  Inner products from $z$ to either side are unchanged, all entries
inside a side are unchanged, and
\[
 \|p_U+p_0\|^2+\|p_W\|^2=\|g_z\|^2.
\]
Consequently the two descendant matrices, embedded back into $V(H)$, add
entrywise to $N$.  In particular, every edge entry is preserved in its unique
descendant.  Repeating this construction terminates precisely at the blocks
of $H$, independently of the order in which cut vertices are split.

Consider such a no-cut descendant $(C,X)$, and write $h=|V(C)|$,
$\ell=|E(C)|$, $q=q(C)$, $d_0=\tr X$, and
$w=\sum_{uv\in E(C)}X_{uv}$.  Equality cannot pass through the strict second
case of the Liu--Tang--Zhang argument recalled after
\eqref{eq:threshold}.  It therefore lies in the first case, and equality in
the final bound forces equality both in the flat edge Cauchy--Schwarz
estimate and at the threshold:
\begin{equation}\label{eq:no-cut-equalities}
 X_{uv}=r\quad(uv\in E(C)),
 \qquad 2(h-1)w=qd_0.
\end{equation}
Here $r>0$: the tight terminal descendant has $T=\lambda q>0$, and hence
$S>0$.  Since
$w=\ell r$, equality in \eqref{eq:flat} is explicit.

The averaged estimate \eqref{eq:averaged} remains valid independently of
which case proves the theorem.  Comparing it with $4S^2=qT$ gives
\begin{equation}\label{eq:d0-lower}
 d_0\geq\frac{h-2}{q-1}T.
\end{equation}
On the other hand, the threshold equality in
\eqref{eq:no-cut-equalities}, together with $T=d_0+2w$, gives
\begin{equation}\label{eq:d0-exact}
 d_0=\frac{h-1}{q+h-1}T.
\end{equation}
Comparing \eqref{eq:d0-lower} and \eqref{eq:d0-exact} and
cross-multiplying yields
\[
 (h-1)(q-1)\geq(h-2)(q+h-1),
 \qquad\text{hence}\qquad q\geq(h-1)^2.
\]
Simplicity supplies the reverse bound
\[
 q=2\ell-h+1\leq h(h-1)-h+1=(h-1)^2.
\]
Thus $\ell=\binom h2$ and $C=K_h$.  This is the crucial complete-block
conclusion; it comes from retaining the averaged estimate even though the
flat estimate is the one used in the first case.

It remains to identify the terminal matrix rather than only its edge
entries.  For $C=K_h$, we have $q=(h-1)^2$, and equality in the flat estimate
gives
\[
 qT=4\binom h2^2r,
 \qquad T=h^2r.
\]
Equation \eqref{eq:d0-exact} gives $d_0=hr$.  Write the diagonal entries of
$X$ as $rb_1,\ldots,rb_h$.  Every $2\times2$ principal minor is
nonnegative, so $b_ib_j\geq1$ for $i\ne j$, while
$\sum_i b_i=h$.  These conditions force $b_i=1$ for every $i$: if some
$b_i<1$, every other $b_j\geq b_i^{-1}$ and the sum exceeds $h$; if no
$b_i<1$, the sum condition is decisive.  Consequently $X=rJ_h$.  Since
$T=\lambda q$, we obtain
\begin{equation}\label{eq:complete-atom}
 X=\lambda\frac{(h-1)^2}{h^2}J_h
   =\lambda\frac{q(C)}{|V(C)|^2}J_{V(C)}.
\end{equation}

For a bridge $C=K_2$, the base case is the same conclusion.  Indeed, if
$X=\bigl(\begin{smallmatrix}a&c\\c&b\end{smallmatrix}\bigr)$ is tight,
then
\[
 4c=a+b+2c,
 \]
whereas positive semidefiniteness gives
$2c\leq2\sqrt{ab}\leq a+b$.  Equality follows throughout, so
$a=b=c$, and $X=(\lambda/4)J_2$ because $q(K_2)=1$ and $T=\lambda$.

The entrywise cut identity proved above recombines the terminal matrices by
adding their embedded block contributions; the only overlap is on cut-vertex
diagonal entries, where the contributions add to the parent diagonal.
Induction on the number of cut splits now gives \eqref{eq:block-atoms}.
Finally, an edge belongs to a unique block.
For a bridge,
$\sqrt{q(C)}/|V(C)|=1/2=q(C)/(2|E(C)|)$; for $C=K_h$, both quantities equal
$(h-1)/h$.  Taking the square root of its coefficient in
\eqref{eq:block-atoms} proves \eqref{eq:edge-root}.
\end{proof}

We apply the proposition to the negative spectral part.  This also explains
why signs and absolute values cannot be omitted from the initial reduction.

\begin{corollary}\label{cor:block-graph}
Suppose $G$ is connected, has $n\geq2$ vertices and $m$ edges, and satisfies
$s^+(G)=n-1$.  Then every block of $G$ is complete.  More precisely, if an
edge $uv$ lies in a block $C$, then
\begin{equation}\label{eq:B-edge}
 B_{uv}=-\frac{q(C)}{2|E(C)|}.
\end{equation}
\end{corollary}

\begin{proof}
Put $q=2m-n+1$.  By \eqref{eq:sum-energy}, $s^-(G)=q$.  The Schur product
theorem~\cite[Theorem~7.5.3]{HornJohnson} makes $M=B\circ B$ doubly
nonnegative.  Moreover
\[
 \tr(AB)=\tr((P-B)B)=-\tr B^2=-s^-(G),
\]
and $A$ has two symmetric unit entries for each edge.  Therefore
\[
 \frac{s^-(G)}2=-\sum_{uv\in E(G)}B_{uv}
 \leq\sum_{uv\in E(G)}|B_{uv}|.
\]
Theorem~\ref{thm:ltz}, applied to $B\circ B$, now gives the requested
starting chain
\begin{equation}\label{eq:starting-chain}
 s^-(G)^2
 \leq4\left(\sum_{uv\in E(G)}|B_{uv}|\right)^2
 \leq q\,\one^T(B\circ B)\one
 =q\,s^-(G).
\end{equation}
The first inequality follows exactly from the displayed trace identity and
the triangle inequality on the edge entries.  Since $s^-(G)=q>0$, both
inequalities in \eqref{eq:starting-chain} are equalities.  Equality in the
triangle inequality gives
\begin{equation}\label{eq:sign}
 B_{uv}\leq0\qquad(uv\in E(G)),
\end{equation}
and $B\circ B$ is an equality matrix for Theorem~\ref{thm:ltz}.  Its
normalization in Proposition~\ref{prop:trace} is
$\lambda=T(B\circ B)/q=s^-(G)/q=1$.  The proposition makes every block a
bridge or a complete graph and, because folding preserves edge entries,
gives $|B_{uv}|=q(C)/(2|E(C)|)$.  Equation \eqref{eq:sign} proves
\eqref{eq:B-edge}.
\end{proof}

Thus a putative equality graph is a \emph{block graph}: its blocks are
complete.  Notice that this conclusion does not assert, falsely, that a leaf
partition turns all cyclic blocks into triangles.  Blocks of arbitrary order
remain possible at this stage and will all be excluded at once.

\section{Strict positivity between siblings}\label{sec:siblings}

For a symmetric matrix $K$, let $K_+$ denote its positive spectral part.
The next lemma is the strict input for the final variational argument.

\begin{lemma}[Sibling positivity]\label{lem:siblings}
Let $T$ be a tree.  If distinct vertices $u,v$ have a common neighbor in
$T$, then
\[
 (A(T)_+)_{uv}>0.
\]
\end{lemma}

\begin{proof}
Choose a bipartition $L\sqcup R$ of $T$ with $u,v\in L$, and write
\begin{equation}\label{eq:bip-block}
 A(T)=\begin{pmatrix}0&X\\X^T&0\end{pmatrix},
 \qquad Q=XX^T.
\end{equation}
For $r\in R$, let $j_r\in\R^L$ be the indicator of $N_T(r)$.  For every
$a>0$,
\begin{equation}\label{eq:Q-hypergraph}
 Q+aI=aI+\sum_{r\in R}j_rj_r^T.
\end{equation}
We prove
\begin{equation}\label{eq:inverse-negative}
 (Q+aI)^{-1}_{uv}<0.
\end{equation}

Regard $T$ itself as the incidence graph between the variable vertices $L$
and the neighborhood hyperedges $N_T(r)$, $r\in R$.  It is a tree.  Root it
at a common neighbor $r_0$ of $u$ and $v$, and eliminate branches in
postorder.  We spell out the Schur calculation, including singleton
neighborhoods.

Consider a nonroot neighborhood node $r$ whose parent variable is $p\in L$.
After all branches below it have been eliminated, its child variables
$C=N_T(r)\setminus\{p\}$ carry a positive diagonal matrix
$D_C$: positivity starts with the term $aI$, and the calculation below shows
that every eliminated child neighborhood adds a positive number to its
articulation diagonal.  On $\{p\}\cup C$, the contribution of $r$ together
with those effective child diagonals has the form $J+D_C$ on the child
coordinates, with the $p$ diagonal of $J$ kept at this step.  Eliminating
$C$ changes the $p$ diagonal by
\begin{align}
 1-\one_C^T(D_C+J_C)^{-1}\one_C
 &=1-\frac{\one_C^TD_C^{-1}\one_C}
 {1+\one_C^TD_C^{-1}\one_C} \notag\\
 &=\frac{1}{1+\one_C^TD_C^{-1}\one_C}>0,
 \label{eq:branch-response}
\end{align}
where Sherman--Morrison was used.  If $C=\varnothing$, so that $r$ is a
singleton leaf neighborhood, the value in \eqref{eq:branch-response} is
understood as $1$; thus such leaves are retained rather than silently
discarded.  This proves inductively both the asserted diagonal form and its
strict positivity.

After every branch below $r_0$ is eliminated, the Schur complement on the
root neighborhood $N_T(r_0)$ is
\[
 D_0+J,
 \qquad D_0=\operatorname{diag}(d_x:x\in N_T(r_0)),
 \qquad d_x>0.
\]
Schur complementation preserves the corresponding entries of the inverse.
Another application of Sherman--Morrison gives, for distinct
$x,y\in N_T(r_0)$,
\begin{equation}\label{eq:root-inverse}
 (D_0+J)^{-1}_{xy}
 =-\frac{d_x^{-1}d_y^{-1}}
 {1+\sum_{z\in N_T(r_0)}d_z^{-1}}<0.
\end{equation}
Taking $x=u$ and $y=v$ proves \eqref{eq:inverse-negative}.  This argument is
valid without any rank assumption on $X$ because $a>0$ makes every matrix
being eliminated positive definite.

The scalar identity
\[
 \sqrt t=\frac1\pi\int_0^\infty
 \frac{t}{t+a}\,a^{-1/2}\,da\qquad(t\geq0)
\]
and spectral functional calculus~\cite[Chapter~V]{Bhatia} give, for
$u\ne v$,
\begin{equation}\label{eq:resolvent}
 (Q^{1/2})_{uv}
 =-\frac1\pi\int_0^\infty
 a^{1/2}(Q+aI)^{-1}_{uv}\,da.
\end{equation}
The integral converges entrywise.  As $a\downarrow0$,
$(Q+aI)^{-1}=a^{-1}\Pi_{\ker Q}+O(1)$, so after multiplication by
$a^{1/2}$ the possible singularity is $O(a^{-1/2})$.  As $a\to\infty$, an
off-diagonal resolvent entry is $O(a^{-2})$, and the integrand is
$O(a^{-3/2})$.  By
\eqref{eq:inverse-negative}, the integrand after the minus sign is strictly
positive for every $a>0$.  Hence $(Q^{1/2})_{uv}>0$.

Finally, standard functional calculus applied to \eqref{eq:bip-block} gives
\[
 |A(T)|=
 \begin{pmatrix}Q^{1/2}&0\\0&(X^TX)^{1/2}\end{pmatrix},
 \qquad
 A(T)_+=\frac12
 \begin{pmatrix}Q^{1/2}&X\\X^T&(X^TX)^{1/2}\end{pmatrix}.
\]
Thus $(A(T)_+)_{uv}=\tfrac12(Q^{1/2})_{uv}>0$, as claimed.
\end{proof}

\section{The variational finish}

We use the following standard identity, included for completeness.

\begin{lemma}\label{lem:variational}
For every real symmetric matrix $K$,
\begin{equation}\label{eq:variational}
 \tr(K_+^2)=\max_{Y\succeq0}
 \bigl\{2\tr(KY)-\tr(Y^2)\bigr\}.
\end{equation}
\end{lemma}

\begin{proof}
Since $K=K_+-K_-$ and $K_-,Y\succeq0$,
$\tr(K_-Y)\geq0$.  Therefore
\[
 2\tr(KY)-\tr(Y^2)
 \leq2\tr(K_+Y)-\tr(Y^2)
 =\tr(K_+^2)-\|Y-K_+\|_F^2.
\]
Equality is attained at $Y=K_+$.
\end{proof}

\begin{proposition}\label{prop:block-strict}
Every connected block graph $G$ on $n$ vertices that is not a tree satisfies
\[
 s^+(G)>n-1.
\]
\end{proposition}

\begin{proof}
Consider the augmented incidence graph whose nodes are all vertices and all
blocks of $G$, with a vertex joined to each block containing it.  For a block
graph this incidence graph is a tree.  Root it at an arbitrary graph vertex.
In each complete block $C$, let $p_C$ be its unique vertex nearest the root.
Retain only the star edges $p_Cx$ for
$x\in V(C)\setminus\{p_C\}$.  For a bridge this retains its sole edge.  The
union of all retained stars is connected and has
\[
 \sum_C(|V(C)|-1)=n-1
\]
edges, so it is a spanning tree $T$ of $G$.

Every deleted edge $uv$ belongs to a complete block and joins two distinct
vertices of $V(C)\setminus\{p_C\}$.  Thus $u$ and $v$ have the common
neighbor $p_C$ in $T$.  Write
\[
 A(G)=A(T)+M,
 \qquad P=A(T)_+,
\]
where $M$ is the adjacency matrix of the deleted edges.  Testing $Y=P$ in
\eqref{eq:variational} gives
\begin{align}
 s^+(G)
 &\geq2\tr((A(T)+M)P)-\tr(P^2)\notag\\
 &=s^+(T)+2\tr(MP)\notag\\
 &=n-1+4\sum_{uv\in E(G)\setminus E(T)}P_{uv}.
 \label{eq:strict-variational}
\end{align}
Here $s^+(T)=|E(T)|=n-1$, because the spectrum of the bipartite graph $T$
is symmetric and $s^+(T)+s^-(T)=2|E(T)|$.  Every summand in
\eqref{eq:strict-variational} is strictly positive by
Lemma~\ref{lem:siblings}.  Since a non-tree block graph has at least one
deleted edge, the inequality is strict.
\end{proof}

\begin{proof}[Proof of Theorem~\ref{thm:main}]
First let $n=1$.  The unique connected graph is $K_1$, which is a tree, and
$s^+(K_1)=0=n-1$.

Now let $n\geq2$ and suppose $s^+(G)=n-1$.  Corollary~\ref{cor:block-graph}
shows that every block of $G$ is complete, so $G$ is a block graph.  If it
were not a tree, Proposition~\ref{prop:block-strict} would give
$s^+(G)>n-1$, a contradiction.  Hence $G$ is a tree.

Conversely, if $G$ is a tree, then it is bipartite, so its nonzero adjacency
eigenvalues occur in opposite pairs.  Therefore
\[
 s^+(G)=\frac12\tr A(G)^2=|E(G)|=n-1.
\]
This also covers $K_1$ directly, and completes the proof.
\end{proof}

\begin{thebibliography}{9}

\bibitem{AKMPZ}
S.~Akbari, H.~Kumar, B.~Mohar, S.~Pragada, and S.~Zhang,
\emph{Refinement of a conjecture on positive square energy of graphs},
arXiv:2506.07264, 2025.

\bibitem{Bhatia}
R.~Bhatia,
\emph{Matrix Analysis},
Graduate Texts in Mathematics 169, Springer, New York, 1997.

\bibitem{EFGW}
C.~Elphick, M.~Farber, F.~Goldberg, and P.~Wocjan,
\emph{Conjectured bounds for the sum of squares of positive eigenvalues of a
graph},
Discrete Math. 339 (2016), 2215--2223.

\bibitem{HornJohnson}
R.~A.~Horn and C.~R.~Johnson,
\emph{Matrix Analysis}, 2nd ed.,
Cambridge University Press, Cambridge, 2013.

\bibitem{LTZ}
Y.~Liu, Q.~Tang, and S.~Zhang,
\emph{The positive and negative square-energy conjecture},
arXiv:2607.18031, 2026.

\end{thebibliography}

\end{document}
